% ============================================================
%  Übungsaufgaben Template (my_dhbw Stil, uebung_*.tex)
%  Header "Übungsblatt", grün eingefärbte <details>-Lösungen.
%  Renderer: pdflatex (zweimal)
% ============================================================
\documentclass[11pt, a4paper]{article}

\usepackage[ngerman]{babel}
\usepackage{fontspec}
\setmainfont[
  Path=/usr/share/fonts/TTF/,
  Extension=.ttf,
  BoldFont=DejaVuSans-Bold,
  ItalicFont=DejaVuSans-Oblique,
  BoldItalicFont=DejaVuSans-BoldOblique
]{DejaVuSans}
\setsansfont[
  Path=/usr/share/fonts/TTF/,
  Extension=.ttf,
  BoldFont=DejaVuSans-Bold,
  ItalicFont=DejaVuSans-Oblique,
  BoldItalicFont=DejaVuSans-BoldOblique
]{DejaVuSans}
\setmonofont[Path=/usr/share/fonts/TTF/,Extension=.ttf]{DejaVuSansMono}
\usepackage[top=2cm, bottom=2cm, left=2.4cm, right=2.4cm, footskip=1cm]{geometry}
\usepackage{amsmath, amssymb}
\usepackage{xcolor}
\usepackage{enumitem}
\usepackage{fancyhdr}
\usepackage{lastpage}
\usepackage{tabularx}
\usepackage{booktabs}
\usepackage{microtype}
\usepackage{parskip}

\usepackage{array}
\usepackage{longtable}
\usepackage{ltablex}
\keepXColumns
\usepackage{colortbl}
\usepackage[most,breakable]{tcolorbox}
\usepackage{listings}
\usepackage[normalem]{ulem}
\usepackage{pdflscape}
\usepackage{titlesec}
\usepackage[colorlinks=true, linkcolor=accent, urlcolor=accent]{hyperref}

\renewcommand{\familydefault}{\sfdefault}

% ── Theme ─────────────────────────────────────────────────────
\definecolor{accent}{RGB}{0, 60, 113}
\definecolor{primaryblue}{RGB}{0, 60, 113}
\definecolor{loesung}{RGB}{0, 100, 0}
\definecolor{codebg}{RGB}{245, 245, 245}
\definecolor{figurebg}{RGB}{232, 241, 255}
\definecolor{tablehintbg}{RGB}{240, 255, 240}
\definecolor{solutionbg}{RGB}{240, 252, 240}
\definecolor{rulegray}{RGB}{180, 180, 180}
\definecolor{tableheadbg}{RGB}{0, 60, 113}

% ── Header/Footer ─────────────────────────────────────────────
\pagestyle{fancy}
\fancyhf{}
\renewcommand{\headrulewidth}{0.4pt}
\fancyhead[L]{\footnotesize\color{accent}\nichttitle}
\fancyhead[R]{\footnotesize\color{accent}\nichtsubtitle}
\fancyfoot[R]{\footnotesize\color{accent}Blatt \thepage\,/\,\pageref{LastPage}}

% ── Typografie ────────────────────────────────────────────────
\titleformat{\section}
    {\color{accent}\large\bfseries}{}{0pt}{}
    [\vspace{-4pt}\color{accent}\rule{\linewidth}{1.5pt}]
\titleformat{\subsection}
    {\normalsize\bfseries\color{accent!85!black}}{}{0pt}{}{}
\titleformat{\subsubsection}
    {\small\bfseries\color{accent!70!black}}{}{0pt}{}{}
\titlespacing*{\section}{0pt}{10pt}{3pt}
\titlespacing*{\subsection}{0pt}{6pt}{2pt}
\titlespacing*{\subsubsection}{0pt}{4pt}{1pt}

\setlength{\parindent}{0pt}
\setlength{\parskip}{6pt}
\renewcommand{\arraystretch}{1.2}
\setlist[itemize]{noitemsep, topsep=3pt, parsep=0pt, leftmargin=1.4em}
\setlist[enumerate]{noitemsep, topsep=3pt, parsep=0pt, leftmargin=1.6em}

% ── Boxen: solutionbox = Lösungsbox in grün ──────────────────
\newtcolorbox{figurebox}[1]{
  colback=figurebg, colframe=accent!50,
  title={Abbildung: #1}, fonttitle=\small\bfseries,
  breakable, before skip=6pt, after skip=6pt
}
\newtcolorbox{tablehintbox}[1]{
  colback=tablehintbg, colframe=green!50!black!50,
  title={Tabelle: #1}, fonttitle=\small\bfseries,
  breakable, before skip=6pt, after skip=6pt
}
\newtcolorbox{solutionbox}[1]{
  enhanced,
  colback=solutionbg, colframe=loesung,
  title={#1}, fonttitle=\small\bfseries,
  coltitle=white, colbacktitle=loesung,
  breakable, before skip=8pt, after skip=8pt
}

\lstset{
  basicstyle=\ttfamily\small,
  backgroundcolor=\color{codebg},
  breaklines=true,
  frame=single, framerule=0pt,
  xleftmargin=4pt, xrightmargin=4pt,
  aboveskip=6pt, belowskip=6pt,
  keepspaces=true
}

\providecommand{\nichttitle}{Übungsblatt}
\providecommand{\nichtsubtitle}{}

\renewcommand{\nichttitle}{Betriebssysteme}
\renewcommand{\nichtsubtitle}{Übungsblatt 12}


\begin{document}

\begin{center}
{\Large\bfseries\color{accent}\nichttitle}
\ifx\nichtsubtitle\empty\else
  \\[2pt]{\normalsize\color{accent!80!black}\nichtsubtitle}
\fi
\\[2pt]
\color{accent}\rule{0.4\linewidth}{0.8pt}\color{black}
\end{center}
\vspace{4pt}

% ── BODY ──────────────────────────────────────────────────────

\section{Übungsblatt 12 — Virtueller Speicher}


\noindent\textcolor{rulegray}{\rule{\linewidth}{0.4pt}}


\paragraph{Aufgabe 1 — Virtueller Adressraum \& Konsequenzen \textit{(12 Punkte)}}\mbox{}\\[2pt]

a) Erkläre den Begriff \textbf{virtueller Adressraum} in eigenen Worten und nenne, in welchem Verhältnis er zum physischen RAM steht. \textit{(4 P)}


b) Im Kapitel heißt es: \textit{„Speicherbedarf darf > RAM sein."} Begründe, warum das überhaupt möglich ist, und nenne den Speicherbereich, der diese Eigenschaft technisch ermöglicht. \textit{(4 P)}


c) Ein Entwickler argumentiert: „Solange ich genug RAM kaufe, brauche ich keinen virtuellen Speicher — der Adressraum ist dann ohnehin physisch vorhanden." Widerlege diese Aussage mit \textbf{zwei} Vorteilen aus dem Kapitel, die unabhängig von der RAM-Größe gelten. \textit{(4 P)}


\textbf{Lösung:}


a) Der virtuelle Adressraum ist der Speicher, den ein Prozess vom Betriebssystem zugewiesen bekommt — er existiert nur logisch / „virtuell". Jeder Prozess sieht \textbf{seinen eigenen} Adressraum, unabhängig von den anderen Prozessen. Nur der aktuell benötigte Teil liegt im RAM, der Rest auf dem Sekundärspeicher. Der virtuelle Adressraum ist also eine vom physischen RAM entkoppelte Sicht.


b) Möglich ist das, weil nur die \textit{aktuell benötigten} Teile des Adressraums im RAM liegen müssen. Alle übrigen Seiten werden auf der \textbf{Paging Area (Schattenspeicher)} auf der Festplatte gehalten und bei Bedarf nachgeladen. Der adressierbare Speicher kann dadurch deutlich größer als der physische RAM sein.


c) Auch bei beliebig viel RAM bleiben aus dem Kapitel:

\begin{itemize}
  \item \textbf{Externe Fragmentierung untergeordnet}: Da Seiten und Seitenrahmen gleich groß sind, entsteht kein Verschnitt zwischen frei werdenden Bereichen.
  \item \textbf{Eigener Adressraum pro Prozess}: Prozesse stören sich nicht gegenseitig; der Schutz ist eine Eigenschaft des Adressraum-Konzepts, nicht der RAM-Menge.
\end{itemize}

(Alternativ akzeptabel: „mehr Programme gleichzeitig im RAM → CPU-Auslastung ↑" als organisatorischer Vorteil.)



\noindent\textcolor{rulegray}{\rule{\linewidth}{0.4pt}}


\paragraph{Aufgabe 2 — Seiten, Seitenrahmen, Paging Area berechnen \textit{(18 Punkte)}}\mbox{}\\[2pt]

Gegeben sei ein System mit folgenden Eckdaten:


\begin{itemize}
  \item physischer RAM: \textbf{256 MiB}
  \item Seitengröße / Seitenrahmengröße: \textbf{4 KiB}
  \item ein Prozess P mit virtuellem Adressraum \textbf{1 GiB}, davon zur Zeit \textbf{48 MiB} „aktiv" (im RAM)
  \item ein zweiter Prozess Q mit virtuellem Adressraum \textbf{512 MiB}, davon \textbf{16 MiB} im RAM
\end{itemize}

Beantworte:


a) Wie viele Seitenrahmen umfasst der RAM insgesamt? \textit{(3 P)}


b) Wie viele Seiten umfassen die virtuellen Adressräume von P bzw. Q? \textit{(4 P)}


c) Wie viele Seitenrahmen sind durch P und Q zusammen aktuell belegt? Wie viele Seitenrahmen bleiben für das BS / weitere Prozesse übrig? \textit{(4 P)}


d) Wie groß muss die \textbf{Paging Area} mindestens sein, wenn die \textit{vollständigen} virtuellen Adressräume beider Prozesse jederzeit auslagerbar sein sollen (Worst Case, alles ausgelagert)? \textit{(4 P)}


e) Begründe kurz, warum man in der Praxis selten mit einer Seitengröße von \textbf{1 KiB} arbeitet, obwohl das Kapitel sie als möglich nennt. \textit{(3 P)}


\textbf{Lösung:}


a) RAM = 256 MiB = 256 · 1024 KiB = 262 144 KiB.

Seitenrahmen = 262 144 KiB / 4 KiB = \textbf{65 536 Seitenrahmen} (= 2¹⁶).


b) P: 1 GiB = 1 048 576 KiB → 1 048 576 / 4 = \textbf{262 144 Seiten}.

Q: 512 MiB = 524 288 KiB → 524 288 / 4 = \textbf{131 072 Seiten}.


c) Aktiv im RAM:

\begin{itemize}
  \item P: 48 MiB / 4 KiB = 49 152 / 4 = \textbf{12 288 Seitenrahmen}
  \item Q: 16 MiB / 4 KiB = 16 384 / 4 = \textbf{4 096 Seitenrahmen}
  \item zusammen: 12 288 + 4 096 = \textbf{16 384 Seitenrahmen} belegt
  \item frei: 65 536 − 16 384 = \textbf{49 152 Seitenrahmen} (≙ 192 MiB).
\end{itemize}

d) Worst Case ist „alles ausgelagert", also der gesamte virtuelle Adressraum beider Prozesse:

1 GiB + 512 MiB = \textbf{1,5 GiB = 1 536 MiB}.

Die Paging Area muss also mindestens 1 536 MiB groß sein. (Hinweis: Da der virtuelle Adressraum > RAM ist, reicht es nicht, sich an der RAM-Größe zu orientieren — der Schattenspeicher dimensioniert sich am virtuellen Bedarf.)


e) Bei 1 KiB-Seiten gibt es viermal so viele Seiten/Seitenrahmen wie bei 4 KiB. Verwaltungsaufwand (Tabellen, Buchführung pro Seite) und I/O-Overhead beim Auslagern steigen entsprechend, der Nutzeffekt feinerer Granularität ist dagegen meist gering — daher sind 4–16 KiB praktischer Standard.



\noindent\textcolor{rulegray}{\rule{\linewidth}{0.4pt}}


\paragraph{Aufgabe 3 — Transfer: Multiprogramming-Effekt \textit{(15 Punkte)}}\mbox{}\\[2pt]

Im Kapitel wird als Vorteil 2 genannt: \textit{„Mehr Programme im RAM → CPU-Auslastung ↑"}.


Betrachte folgendes Szenario (im Kapitel nicht beschrieben, aber mit Kapitelwissen lösbar):


\begin{quote}
Ein Server hat 2 GiB RAM. Ein Programm A benötigt durchschnittlich \textbf{1,6 GiB} und verbringt 80 \% seiner Zeit blockiert auf I/O. Auf dem System soll auch ein Programm B (durchschnittlich 1,2 GiB, ebenfalls 70 \% I/O-Wartezeit) laufen.
\end{quote}

a) Erkläre, warum das \textbf{ohne} virtuellen Speicher in diesem Szenario nicht funktionieren würde. \textit{(4 P)}


b) Erkläre, warum es \textbf{mit} virtuellem Speicher funktioniert, obwohl 1,6 GiB + 1,2 GiB > 2 GiB sind. Verwende die Begriffe \textit{Seite}, \textit{Seitenrahmen} und \textit{Paging Area}. \textit{(6 P)}


c) Begründe mithilfe der I/O-Zeiten, warum dadurch die CPU-Auslastung tatsächlich steigt — und nicht nur der Speicherverbrauch. \textit{(5 P)}


\textbf{Lösung:}


a) Ohne virtuellen Speicher müsste \textbf{der gesamte Adressraum eines Programms} physisch im RAM liegen. Da 1,6 GiB + 1,2 GiB = 2,8 GiB > 2 GiB, könnte höchstens eines der Programme zur selben Zeit gehalten werden; das andere müsste komplett geladen/entladen werden (Swapping ganzer Prozesse). Gleichzeitiger Lauf wäre nur durch ständiges vollständiges Umladen möglich — sehr teuer.


b) Mit virtuellem Speicher liegt nicht der ganze Adressraum, sondern nur die \textbf{aktuell benötigten Seiten} in den RAM-\textbf{Seitenrahmen}. Inaktive Seiten beider Programme bleiben in der \textbf{Paging Area} auf der Festplatte. RAM enthält also eine Mischung „heißer" Seiten von A und B; bei Bedarf wird eine Seite aus der Paging Area nachgeladen, eine andere ggf. ausgelagert. So passen die \textit{aktiven Arbeitsmengen} beider Programme zusammen in 2 GiB, obwohl die Summe der virtuellen Adressräume größer ist.


c) Wenn nur A liefe, wäre die CPU 80 \% der Zeit untätig (Wartet auf I/O). Mit B im RAM kann die CPU in A's Wartephasen B bedienen (und umgekehrt). Da beide Programme oft blockieren, liegt die Wahrscheinlichkeit, dass \textit{beide gleichzeitig} warten, deutlich niedriger (≈ 0,8 · 0,7 = 56 \% statt 80 \%). Die CPU-Auslastung steigt entsprechend von 20 \% auf rund 44 \%. Genau das meint Vorteil 2.



\noindent\textcolor{rulegray}{\rule{\linewidth}{0.4pt}}


\paragraph{Aufgabe 4 — Analyse / Entscheidung: Swapping vs. Paging \textit{(20 Punkte)}}\mbox{}\\[2pt]

Gegeben sind drei Systeme:



{\small
\begin{tabularx}{\linewidth}{XXXX}
\hline
\rowcolor{tableheadbg}
{\color{white}\textbf{System}} & {\color{white}\textbf{RAM}} & {\color{white}\textbf{MMU vorhanden?}} & {\color{white}\textbf{typische Last}} \\[2pt]
\hline
\endhead
\hline
\endfoot
\textbf{S1} Mikrocontroller (Heizungssteuerung) & 64 KiB & nein & 1 Steuerprozess, 2 KiB Code \\
\textbf{S2} Bürorechner & 16 GiB & ja & Browser, IDE, Mailclient gleichzeitig \\
\textbf{S3} Embedded-Linux-Kasten an einer Maschine & 128 MiB & ja & 1 Hauptprozess (90 MiB) + selten genutzte Diagnose-App (60 MiB) \\
\end{tabularx}
}


a) Wähle für jedes System begründet \textbf{Swapping} \textit{oder} \textbf{Paging} (oder „keines von beiden nötig"). Stütze dich auf die im Kapitel genannten Unterscheidungsmerkmale \textbf{Granularität} und \textbf{Einsatz (mit/ohne MMU)}. \textit{(9 P)}


b) Erkläre für \textbf{S2} explizit, welche der vier Vorteile aus dem Kapitel hier wie zum Tragen kommen (jeden Vorteil benennen und auf das Szenario beziehen). \textit{(8 P)}


c) Für \textbf{S3} entscheidet sich der Hersteller dennoch für Swapping. Nenne einen plausiblen Trade-off, der diese Entscheidung trotz vorhandener MMU rechtfertigen kann. \textit{(3 P)}


\textbf{Lösung:}


a)

\begin{itemize}
  \item \textbf{S1} — \textit{keines von beiden nötig (alternativ: Swapping bei zusätzlichen Prozessen)}. Ohne MMU ist Paging laut Kapitel nicht einsetzbar (Paging setzt MMU voraus). Da nur ein 2-KiB-Prozess existiert und 64 KiB RAM ausreichen, ist faktisch gar kein Auslagern nötig. \textit{Müsste} doch ausgelagert werden, käme nur \textbf{Swapping} in Frage (Granularität: ganzer Prozess).
  \item \textbf{S2} — \textit{Paging}. MMU vorhanden, viele Prozesse gleichzeitig, summierter Speicherbedarf > RAM realistisch. Feine Granularität (Teile eines Prozesses) erlaubt es, dass alle Programme nebenher leben — genau das Einsatzgebiet von Paging.
  \item \textbf{S3} — \textit{Paging (Standardentscheidung)}. MMU ist da; Hauptprozess + Diagnose-App können sich Seiten teilen, RAM (128 MiB) wird sonst knapp, sobald Diagnose läuft.
\end{itemize}

b) Für S2:

\begin{enumerate}
  \item \textbf{Programme > RAM möglich}: Einzelne Anwendungen (z. B. IDE mit großem Projekt) können größeren virtuellen Adressraum nutzen, als physisch verfügbar wäre.
  \item \textbf{Mehr Programme im RAM → CPU-Auslastung ↑}: Browser, IDE, Mailclient blockieren oft auf I/O (Netzwerk, Disk); die CPU kann in den Wartephasen anderen Prozessen dienen.
  \item \textbf{Ansprechbarer Speicher > physisch}: Jeder Prozess kann einen logisch großen Adressraum adressieren, unabhängig von 16 GiB.
  \item \textbf{Externe Fragmentierung untergeordnet}: Da alle Seitenrahmen gleich groß sind, fragmentiert der RAM beim ständigen Starten/Beenden von Programmen nicht.
\end{enumerate}

c) Plausibler Trade-off: Die Diagnose-App läuft \textit{selten} und ihr Code ist i. d. R. \textbf{vollständig} zu laden (wenig Lokalität); Swapping eines kompletten Prozesses ist dann einfacher zu implementieren und braucht keine Seitenverwaltung. Außerdem ist der Speicheraufwand für Seitentabellen / Buchführung in 128 MiB nicht vernachlässigbar — Swapping spart diesen Verwaltungs-Overhead.



\noindent\textcolor{rulegray}{\rule{\linewidth}{0.4pt}}


\paragraph{Aufgabe 5 — Fehler finden in einer Systembeschreibung \textit{(10 Punkte)}}\mbox{}\\[2pt]

Ein Praktikant beschreibt die Speicherorganisation seines Systems wie folgt:


\begin{quote}
\textit{„Wir haben 8 GiB RAM und arbeiten mit Paging. Eine }\textit{Seite}\textit{ ist bei uns 4 KiB groß, ein }\textit{Seitenrahmen}\textit{ 8 KiB — das gibt uns mehr Flexibilität. Die }\textit{Paging Area}\textit{ liegt direkt im RAM, damit das Auslagern schnell geht. Da unser Mikrocontroller keine MMU hat, sparen wir uns Swapping komplett."}
\end{quote}

Identifiziere \textbf{alle} fachlichen Fehler in dieser Aussage und korrigiere sie jeweils mit Begründung aus dem Kapitel.


\textbf{Lösung:}


Drei (bzw. vier) Fehler:


\begin{enumerate}
  \item \textbf{„Seite 4 KiB, Seitenrahmen 8 KiB"} — falsch. Laut Kapitel sind Seitenrahmen Blöcke des RAM und \textbf{alle gleich groß}, und zwar genauso groß wie die Seiten — nur dann passt eine Seite genau in einen Rahmen. Korrektur: Seite und Seitenrahmen müssen identische Größe haben (z. B. beide 4 KiB).
\end{enumerate}

\begin{enumerate}
  \item \textbf{„Paging Area liegt direkt im RAM"} — falsch. Die Paging Area (Schattenspeicher) ist per Definition ein \textbf{Festplattenbereich} für ausgelagerte Seiten. Im RAM zu liegen widerspricht ihrem Zweck — das Auslagern soll ja gerade RAM freigeben.
\end{enumerate}

\begin{enumerate}
  \item \textbf{„Mikrocontroller hat keine MMU → wir nutzen Paging"} — falsch. Paging setzt laut Kapitel-Tabelle eine \textbf{MMU voraus}. Ohne MMU bleibt nur Swapping (ganzer Prozess) — genau das Gegenteil der Aussage des Praktikanten („Swapping sparen wir uns").
\end{enumerate}

\begin{enumerate}
  \item (Folgefehler) Die Aussage „wir arbeiten mit Paging" ist mit „kein MMU" insgesamt unvereinbar — entweder MMU + Paging, oder ohne MMU nur Swapping.
\end{enumerate}

Korrekte Variante: \textit{„Seite und Seitenrahmen sind beide 4 KiB groß. Die Paging Area liegt auf der Festplatte. Ohne MMU ist nur Swapping (ganzer Prozess) möglich, kein Paging."}





% ──────────────────────────────────────────────────────────────

\end{document}
